\documentclass[12pt]{article}
\usepackage{lmodern}
\usepackage{geometry}
\usepackage{enumitem}
\setlist[description]{leftmargin=1cm,labelindent=1cm}
\usepackage{amsmath}
\usepackage{amsthm}
\usepackage{amssymb}
\usepackage{mathtools}
\usepackage[hidelinks]{hyperref}
\usepackage[portuguese]{babel}

\theoremstyle{plain}
\newtheorem{thm}{Teorema}[section]
\newtheorem{prop}[thm]{Proposição}
\newtheorem{lem}[thm]{Lema}
\newtheorem{cor}[thm]{Corolário}
\newtheorem{res}[thm]{Resultado}

\theoremstyle{definition}
\newtheorem{defn}[thm]{Definição}
\newtheorem{exmp}[thm]{Exemplo}
\newtheorem{prblm}[thm]{Problema}
\newtheorem{prcp}[thm]{Princípio}

\theoremstyle{remark}
\newtheorem*{rem}{Nota}

\geometry{
 a4paper,
 left=3cm,
 top=2cm,
 right=2cm,
 bottom=2cm
}

\begin{document}
\begin{center}
    \textbf{\large UNIVERSIDADE FEDERAL DO RIO GRANDE DO SUL} \\
    \textbf{\large INSTITUTO DE MATEMÁTICA E ESTATÍSTICA} \\
    \vspace{3cm}
    \textbf{\Large MATEMÁTICA DISCRETA I} \\
    \vspace{1cm}
    \textbf{\large LISTA DE DEFINIÇÕES E RESULTADOS} \\
    \vspace{3cm}

    \textbf{LEONARDO SUAREZ}

    \vfill
    \textbf{\large Porto Alegre, RS} \\
    \textbf{\large \today}
\end{center}

\newpage
\tableofcontents
\newpage

\section{Métodos de Contagem}

\subsection{Princípios aditivo e multiplicativo}

\begin{prcp} (princípio da adição).
    Se existem \(r_1\) objetos distintos no primeiro conjunto, \(r_2\) no segundo conjunto, \ldots, \(r_n\) no n-ésimo conjunto e esses conjuntos são \emph{disjuntos dois a dois}, então o número de maneiras de selecionar um objeto de um dos \(n\) conjuntos é \(r_1 + r_2 + \cdots + r_n\).
\end{prcp}

\begin{prcp} (princípio da multiplicação).
    Se um processo pode ser dividido em \(n\) etapas, onde a primeira etapa pode ser realizada de \(r_1\) maneiras distintas, a segunda etapa pode ser realizada de \(r_2\) maneiras distintas, \ldots, e a n-ésima etapa pode ser realizada de \(r_n\) maneiras distintas, então o número total de maneiras de realizar o processo é \(r_1 \cdot r_2 \cdots r_n\), dado que a escolha de cada etapa é independente das outras.
\end{prcp}

\begin{prcp} (princípio da composição de conjuntos).
    Dado qualquer conjunto \(S\) construído pelo princípio multiplicativo, é preciso saber dizer quais elementos de \(S\) são da primeira parte de \(S\) e quais são da segunda parte.
\end{prcp}

\subsection{Permutações com repetição}

\begin{thm} (permutações com repetição)
    Se existem \(n\) objetos, com \(r_1\) do tipo 1, \(r_2\) do tipo 2, \ldots, \(r_m\) do tipo \(m\), onde \(r_1 + r_2 + \cdots + r_m = n\), então o número de permutações com esses \(n\) objetos é dado por
    \begin{align}
        P(n;r_1,r_2,\ldots,r_m) &= \binom{n}{r_1}\binom{n-r_1}{r_2}\cdots\binom{n-r_1-r_2-\cdots-r_{m-1}}{r_m} \\
        &= \frac{n!}{r_1!r_2!\cdots r_m!}.
    \end{align}
\end{thm}

\subsection{Combinações com repetição}

\begin{thm} (seleções com repetição).
    O número de seleções com repetição de \(r\) objetos dentre \(n\) tipos distintos é \(C(r+n-1,r)\).
\end{thm}

\begin{prop} (formas equivalentes para seleções com repetição).
    \begin{enumerate}
        \item O número de maneiras de selecionar \(r\) objetos com repetição dentre \(n\) tipos distintos.
        \item O número de maneiras de distribuir \(r\) objetos idênticos em \(n\) caixas distintas.
        \item O número de soluções inteiras não negativas da equação
        \begin{equation}
            x_1 + x_2 + \cdots + x_n = r.
        \end{equation}
    \end{enumerate}
\end{prop}

\subsection{Coeficientes binomiais}

\begin{thm} (teorema binomial).
    Para quaisquer inteiros não-negativos \(n\),
    \begin{equation}
        (x+y)^n = \sum_{k=0}^{n} \binom{n}{k} x^k y^{n-k}.
    \end{equation}
\end{thm}

\begin{prop} (identidades binomiais).
    \begin{gather}
        \binom{n}{0}+\binom{n}{1}+\binom{n}{2}+\cdots+\binom{n}{n} = 2^n \\
        \binom{n}{0}+\binom{n+1}{1}+\binom{n+2}{2}+\cdots+\binom{n+r}{r} = \binom{n+r+1}{r+1} \\
        \binom{r}{r}+\binom{r+1}{r}+\binom{r+2}{r}+\cdots+\binom{n}{r} = \binom{n+1}{r+1} \\
        \binom{n}{0}^2 + \binom{n}{1}^2 + \binom{n}{2}^2 + \cdots + \binom{n}{n}^2 = \binom{2n}{n} \\
        \sum_{k=0}^{r} \binom{m}{k}\binom{n}{r-k} = \binom{m+n}{r} \\
        \sum_{k=0}^{m} \binom{m}{k}\binom{n}{r+k} = \binom{m+n}{m+r} \\
        \sum_{k=n-s}^{m-r} \binom{m-k}{r}\binom{n+k}{s} = \binom{m+n+1}{r+s+1} \\
    \end{gather}
\end{prop}

\section{Inclusão-Exclusão}

\subsection{Cardinalidade da união de \(n\) conjuntos}

\begin{thm} (princípio da inclusão-exclusão)
    Sejam \(A_1, A_2, \ldots, A_n\) subconjuntos finitos de um conjunto universal \(U\). Seja \(S_k\) a soma de todas as interseções de \(k\) conjuntos dentre os \(A_i\). então
    \begin{equation}
        \left| \bigcup_{i=1}^{n} A_i \right| = S_1 - S_2 + S_3 - S_4 + \cdots + (-1)^{n+1} S_n.
    \end{equation}
\end{thm}

\subsection{A função \(\varphi\) de Euler}

\begin{defn}
    Chamamos de \emph{função \(\varphi\) de Euler} a função que atribui a cada inteiro positivo \(m\) o número de inteiros positivos menores ou iguais a \(m\) que são coprimos a \(m\).
\end{defn}

\begin{thm}
    O valor de \(\varphi(m)\) é dado por
    \begin{equation}
        \varphi(m) = m \left( 1 - \frac{1}{p_1} \right) \left( 1 - \frac{1}{p_2} \right) \cdots \left( 1 - \frac{1}{p_t} \right),
    \end{equation}
    \noindent onde \(p_1, p_2, \ldots, p_t\) são os fatores primos distintos de \(m\), isto é, \(m = p_1^{k_1} p_2^{k_2} \cdots p_t^{k_t}\).
\end{thm}

\subsection{Permutações caóticas}

\begin{defn}
    Uma permutação de \(a_1, a_2, \ldots, a_n\) é chamada de \emph{caótica} quando nenhum dos elementos \(a_i\) está em sua posição natural \(i\). O número de permutações caóticas (ou de \emph{desarranjos}) de \(n\) objetos é denotado por \(D_n\).
\end{defn}

\begin{res}
    O número de permutações caóticas de \(n\) objetos é dado por
    \begin{equation}
        D_n = n! \left( 1 - \frac{1}{1!} + \frac{1}{2!} - \frac{1}{3!} + \cdots + \frac{(-1)^n}{n!} \right).
    \end{equation}
\end{res}

\begin{thm}
    Seja \(\vert A \vert = n\) e \(\vert B \vert = k\). Se \(n=k\), para \(n>0\), o número de funções bijetoras \(f: A \to B\) é \(n!\).
\end{thm}

\begin{thm}
    Para \(n \leq k\), o número de funções injetoras \(f: A \to B\) é \(k(k-1)(k-2)\cdots(k-n+1) = \displaystyle\frac{k!}{(k-n)!}\).
\end{thm}

\begin{thm}
    Para \(n \geq k\), o número de funções sobrejetoras \(f: A \to B\) é dado por
    \begin{equation}
        T(n,k) = \sum_{i=0}^{k} (-1)^i \binom{k}{i} (k-i)^n.
    \end{equation}
\end{thm}

\section{Funções Geradoras}

\begin{defn}
    Suponha que \(a_r\) seja o número de maneiras de selecionar \(r\) objetos em um certo problema de contagem. Então \(g(x)\) é a \emph{função geradora ordinária} de \(a_r\) se
    \begin{equation}
        g(x) = a_0 + a_1 x + a_2 x^2 + \cdots = \sum_{r=0}^{\infty} a_r x^r.
    \end{equation}
\end{defn}

\subsection{Cálculo de coeficientes de funções geradoras}

\begin{thm} Sendo \(f(x)\) e \(g(x)\) funções geradoras de \(a_r\) e \(b_r\), respectivamente, então
    \begin{enumerate}
        \item[\textup{(a)}] \(Af(x) + Bg(x)\) é a função geradora de \(Aa_r + Bb_r\), onde \(A\) e \(B\) são constantes.
        \item[\textup{(b)}] \(\displaystyle f(x)g(x) = \sum_{n=0}^{\infty} \left( \sum_{k=0}^{n} a_k b_{n-k} \right) x^n\).
        \item[\textup{(c)}] A função geradora para \(a_0+a_1+a_2+\cdots+a_r\) é igual a \((1+x+x^2+\cdots)f(x)\).
        \item[\textup{(d)}] A função geradora para \((ra_r)\) é igual a \(xf'(x)\).
        \item[\textup{(e)}] \(\displaystyle \int f(x) dx = \sum_{n=0}^{\infty} \frac{a_n}{n+1} x^{n+1}\).
    \end{enumerate}
\end{thm}

\begin{thm} (teorema binomial generalizado).
    \begin{equation}
        (1+x)^n=1+nx+\frac{n(n-1)}{2!}x^2+\frac{n(n-1)(n-2)}{3!}x^3+\cdots = \sum_{k=0}^{\infty} \binom{n}{k} x^k,
    \end{equation}
    \noindent onde
    \[
        \binom{n}{k} = \begin{cases}
            n(n-1)(n-2)\cdots(n-k+1)/k! & \text{se } k \geq 1, \\
            1 & \text{se } k=0, \\
        \end{cases}
    \]
\end{thm}

\begin{prop} (expansões polinomiais).
    \begin{enumerate}
        \item[\textup{(a)}] \(\displaystyle\frac{1-x^{m+1}}{1-x} = 1 + x + x^2 + \cdots + x^m\)
        \item[\textup{(b)}] \(\displaystyle\frac{1}{1-x} = 1 + x + x^2 + x^3 + \cdots\)
        \item[\textup{(c)}] \((1+x)^n = \displaystyle 1 + \binom{n}{1} x + \binom{n}{2} x^2 + \cdots + \binom{n}{n} x^n\)
        \item[\textup{(d)}] \((1-x^m)^n = \displaystyle 1 - \binom{n}{1} x^m + \binom{n}{2} x^{2m} - \cdots + (-1)^n \binom{n}{n} x^{mn}\)
        \item[\textup{(e)}] \(\displaystyle \frac{1}{(1-x)^n} = \displaystyle 1 + \binom{1+n-1}{n-1} x + \binom{2+n-1}{n-1} x^2 + \cdots = \sum_{k=0}^{\infty} \binom{n+k-1}{n-1} x^k\)
        \item[\textup{(f)}] Se \(h(x)=f(x)g(x)\), onde \(\displaystyle f(x)=a_0 + a_1 x + a_2 x^2 + \cdots = \sum_{j=0}^{\infty} a_j x^j\) e \(\displaystyle g(x)=b_0 + b_1 x + b_2 x^2 + \cdots = \sum_{j=0}^{\infty} b_j x^j\), então
        \begin{equation}
            h(x) = a_0b_0 + (a_0b_1 + a_1b_0)x + (a_0b_2 + a_1b_1 + a_2b_0)x^2 + \cdots = \sum_{j=0}^{\infty} \left( \sum_{k=0}^{j} a_k b_{j-k} \right) x^j.
        \end{equation}
    \end{enumerate}
\end{prop}

\subsection{Função geradora exponencial}

A função geradora exponencial é útil quando estamos interessados em contar objetos onde a ordem dos elementos importa, como em permutações ou arranjos. Quando a ordem é irrelevante, usamos a função geradora ordinária.

\begin{defn}
    A série de potências
    \begin{equation}
        g(x) = a_0 + a_1 \frac{x}{1!} + a_2 \frac{x^2}{2!} + a_3 \frac{x^3}{3!} + \cdots = \sum_{r=0}^{\infty} a_r \frac{x^r}{r!}
    \end{equation}
    \noindent é chamada de \emph{função geradora exponencial} da sequência \((a_r)\).
\end{defn}

\subsection{Partições de um inteiro}

\begin{defn}
    Uma \emph{partição} de um inteiro positivo \(n\) é uma coleção de inteiros positivos cuja soma é \(n\).
\end{defn}

\begin{thm} 
    O número \(p_k(n)\) de partições de \(n\) tendo \(k\) como a maior parte é igual ao número \(q_k(n)\) de partições de \(n\) com exatamente \(k\) partes, isto é, \(p_k(n) = q_k(n)\).
\end{thm}

\begin{cor}
    Seja \(P_k(n)\) o número de partições de \(n\) em partes não maiores que \(k\) e \(Q_k(n)\) o número de partições de \(n\) em no máximo \(k\) partes. Então \(P_k(n) = Q_k(n)\).
\end{cor}

\begin{thm}
    \(F(n) = G(n)\), onde \(F(n)\) é o número de partições de \(n\) em que cada parte aparece pelo menos duas vezes e \(G(n)\) é o número de partições de \(n\) em partes maiores do que \(2\) e tais que nenhuma parte é consecutiva.
\end{thm}

\begin{thm}
    O número de partições autoconjugadas de \(n\) é igual ao número de partições de \(n\) em partes ímpares distintas.
\end{thm}

\begin{thm}
    O número de partições de \(n\) em partes distintas é igual ao número de partições de \(n\) em partes ímpares.
\end{thm}

\section{Relações de Recorrência}

\subsection{Resolução de relações de recorrência lineares homogêneas}

\begin{defn}
    Uma relação de recorrência é dita \emph{homogênea} se todos os termos da relação envolvem algum \(a_k\). Quando um termo adicional envolvendo uma função de \(n\) (ou uma constante) aparece na relação de recorrência, como 
    \begin{equation}
        a_n = c_1 a_{n-1} + \cdots + c_k a_{n-k} + f(n),
    \end{equation}
    \noindent a relação é dita \emph{não-homogênea}.
\end{defn}

\begin{res}
    A solução geral de uma relação de recorrência da forma 
    \begin{equation}
        a_n = c_1 a_{n-1} + c_2 a_{n-2} + \cdots + c_r a_{n-r}
    \end{equation}
    \noindent envolve uma soma de soluções individuais da forma \(a_n=\alpha^n\). Para determinar os valores possíveis de \(\alpha\), substituímos \(a_k=\alpha^k\) na relação de recorrência, o que leva à \emph{equação característica}
    \begin{equation}\label{eq:caracteristica}
        \alpha^r - c_1 \alpha^{r-1} - c_2 \alpha^{r-2} - \cdots - c_r = 0.
    \end{equation}
    Se \(\alpha_1, \alpha_2, \ldots, \alpha_r\) são os \(r\) zeros da equação \eqref{eq:caracteristica}, então a solução geral da relação de recorrência é dada por
    \begin{equation}
        a_n = A_1 \alpha_1^n + A_2 \alpha_2^n + \cdots + A_r \alpha_r^n,
    \end{equation}
    \noindent onde \(A_1, A_2, \ldots, A_r\) são constantes determinadas pelas condições iniciais. Dadas os valores iniciais \(a_0', a_1', \ldots, a_{r-1}'\), podemos encontrar as constantes \(A_1, A_2, \ldots, A_r\) resolvendo o sistema de equações
    \begin{equation}
        A_1\alpha_1^k + A_2\alpha_2^k + \cdots + A_r\alpha_r^k = a_k', \quad k=0,1,\ldots,r-1.
    \end{equation}
    As \(r\) equações lineares acima podem ser resolvidas por eliminação gaussiana para determinar as \(r\) constantes \(A_i\). Se a solução da \href{eq:caracteristica}{equação característica} tem raízes múltiplas, digamos \(\alpha_1\) é uma raiz de multiplicidade \(m\), então os termos correspondentes na solução geral são
    \begin{equation}
        A_1 \alpha_1^n + A_2 n \alpha_1^n + A_3 n^2 \alpha_1^n + \cdots + A_m n^{m-1} \alpha_1^n.
    \end{equation}
\end{res}

\subsection{Resolução de relações de recorrência lineares não-homogêneas}

\begin{res}
    A solução geral de uma relação de recorrência não-homogênea da forma
    \begin{equation}\label{eq:nao-homogenea}
        a_n = c_1 a_{n-1} + c_2 a_{n-2} + \cdots + c_r a_{n-r} + f(n)
    \end{equation}
    \noindent é dada por \(a_n = a_n^{(0)} + a^{*}_{n}\), onde \(a_n^{(0)}\) é a solução geral da relação homogênea associada e \(a^{*}_{n}\) é uma solução particular da relação não-homogênea.
\end{res}

\begin{res} (soluções conhecidas para funções \(f(n)\) específicas para \(c\neq 1\)). \vspace{2mm}\newline
    \begin{tabular}[c]{ll}
        \(f(n)\) & \textup{solução particular} \(a^{*}_{n}\) \\
        \hline
        \(d\)\textup{, constante} & \(B\) \\
        \(dn\) & \(B_1n + B_0\) \\
        \(dn^2\) & \(B_2 n^2 + B_1 n + B_0\) \\
        \(d^n\) & \(B d^n\) \\
    \end{tabular}
\end{res}

\section{O princípio da casa dos pombos}

\begin{thm} (princípio da casa dos pombos).
    Se \(n+1\) pombos são colocados em \(n\) gaiolas, então pelo menos uma gaiola contém dois ou mais pombos.
\end{thm}

\begin{thm} (princípio da casa dos pombos generalizado).
    Se \(n\) gaiolas são ocupadas por \(nk+1\) pombos, para algum inteiro \(k \geq 1\), então pelo menos uma gaiola contém \(k+1\) ou mais pombos.
\end{thm}

\begin{thm}\label{eq:pcp}
    Se colocarmos em \(n\) gaiolas \(k\) pombos, então pelo menos uma gaiola deverá conter pelo menos \(\displaystyle \left \lfloor \frac{(k-1)}{n} \right \rfloor + 1\) pombos.
\end{thm}

\begin{proof}
    Suponha, por absurdo, que nenhuma gaiola contém mais que \(\displaystyle \left \lfloor \frac{(k-1)}{n} \right \rfloor\) pombos. Então,
    \begin{equation*}
        n \left \lfloor \frac{(k-1)}{n} \right \rfloor \leq n \frac{(k-1)}{n} = k-1 < k,
    \end{equation*}
    \noindent o que é uma contradição.
\end{proof}

\begin{thm} (Erdős-Szekeres).
    Qualquer sequência de \(n^2+1\) números inteiros distintos contém uma subsequência crescente ou decrescente de \(n+1\) termos.
\end{thm}

\begin{proof}
    Considera a sequência \(a_1, \ldots, a_{n^2+1}\). Seja \(t_i\) o número de termos na maior subsequência crescente começando em \(a_i\). Se algum \(t_i=n+1\), o teorema está demonstrado. 
    Caso contrário, temos que \(1\leq t_i \leq n\) para todo \(i\). Observa que há \(n^2+1\) pombos (\(t_i\)'s) para \(n\) gaiolas (os números \(1, \ldots, n\)). Pelo \eqref{eq:pcp}, pelo menos uma gaiola deverá conter
    \begin{equation*}
        \left \lfloor \frac{n^2+1-1}{n} \right \rfloor + 1 = n+1
    \end{equation*}
    \noindent pombos, isto é, há pelo menos \(n+1\) \(t_i\)'s iguais. Para mostrar que os \(a_i\)'s correspondentes formam uma subsequência decrescente, suponha, por absurdo, que \(a_i < a_j\) (\(a_i \neq a_j\) pois os termos são diferentes) para \(i < j\).
    Logo, encontramos uma subsequência crescente de \(t_{j}+1\) termos que começa em \(a_i\), o que é uma contradição.
\end{proof}

\section{Noções sobre Grafos}

\begin{defn}
    Um \emph{grafo} G é um tripla que consiste em um \emph{conjunto de vértices} \(V(G)\), um \emph{conjunto de arcos} \(E(G)\), e uma relação que associa a cada arco dois vértices (não necessariamente distintos), chamados de \emph{extremidades} do arco.
\end{defn}

\begin{defn}
    Um \emph{laço} é um arco cujas extremidades são iguais. \emph{Arcos múltiplos} são arcos que possuem o mesmo par de extremidades. 
    
    Um \emph{grafo simples} é um grafo que não possui laços nem arcos múltiplos. Especificamos um grafo simples pelo seus conjunto de vértices e conjunto de arcos, sendo o conjunto de arcos uma coleção de pares não ordenados de vértices. 
    
    Quando \(u\) e \(v\) são as extremidades de um arco, eles são \emph{adjacentes} e \emph{vizinhos}, escrevemos, então, \(u \leftrightarrow v\).
\end{defn}

\begin{defn}
    O \emph{complemento} \(\overline{G}\) de um grafo simples \(G\) é o grafo simples com conjunto de vértices \(V(G)\) definido por \(uv \in E(\overline{G})\) sse \(uv \notin E(G)\).
    Um \emph{clique} em um grafo é um conjunto de vértices adjacentes dois a dois. Um \emph{conjunto independente} (ou \emph{conjunto estável}) num grafo é o conjunto de vértices não adjacentes.
\end{defn}

\begin{defn}
    Um grafo \(G\) é \emph{bipartido} se \(V(G)\) é a união de dois conjuntos independentes disjuntos (possivelmente vazios) chamados de \emph{conjuntos partidos} de \(G\).
\end{defn}

\begin{defn}
    Um \emph{caminho} é um grafo simples cujos vértices podem ser ordenados tal que dois vértices são adjacentes sse eles são consecutivos na lista.
    Um \emph{ciclo} é um grafo com o mesmo número de vértices e de arcos cujos vértices podem ser postos em volta dum círculo tal que dois vértices são adjacentes sse eles aparecem consecutivamente em torno do círculo.
\end{defn}

\begin{defn}
    Um \emph{subgrafo} de um grafo \(G\) é um grafo \(H\) tal que \(V(H) \subset V(G)\) e \(E(H) \subset E(G)\) e a atribuição de arcos a vértices em \(H\) é a mesma que em \(G\).
    Então escrevemos \(H \subset G\) e dizemos que "\(G\) contém \(H\)".
    
    Um grafo \(G\) é \emph{conexo} se cada par de vértices em \(G\) pertence a um caminho; caso contrário, \(G\) é \emph{desconexo}.
\end{defn}

\begin{defn}
    Seja \(G\) um grafo sem laços com \(V(G) = \{v_1, \ldots, v_n\}\) e \(E(G) = \{e_1, \ldots, e_m\}\). A \emph{matriz de adjacência} de \(G\), denotada por \(A(G)\), é a matriz \(n\) por \(n\) em que a entrada \(a_{i,j}\)
    é o número de arcos em \(G\) com extremidades \(\{v_i, v_j\}\). A \emph{matriz de incidência} \(M(G)\) é a matriz \(n\) por \(m\) na qual a entrada \(m_{i,j}\) é \(1\) se \(v_i\) é uma extremidade de \(e_j\) e 0 caso contrário.

    Se o vértice \(v\) é uma extremidade do arco \(e\), então \(v\) e \(e\) são \emph{incidentes}. O \emph{grau} do vértice \(v\) (em um grafo sem laços) é o número de arcos incidentes.
\end{defn}

\subsection{(Multi)grafos eulerianos}

\subsection{Isomorfismo}

\subsection{Planaridade}

\end{document}